\documentclass[11pt,a4paper]{article}

\usepackage[utf8]{inputenc}
\usepackage[T1]{fontenc}
\usepackage[french]{babel}
\usepackage[a4paper,margin=2.1cm]{geometry}
\usepackage{amsmath,amssymb,siunitx}
\usepackage{physics}
\usepackage{enumitem}
\usepackage{tabularx}
\usepackage{array}
\usepackage{xcolor}
\usepackage{lmodern}
\usepackage{fancyhdr}
\usepackage{lastpage}
\usepackage{hyperref}
\usepackage{multicol}
\usepackage{titlesec}
\usepackage{amsfonts}
\usepackage{tcolorbox}

\hypersetup{
    colorlinks=true,
    linkcolor=black,
    urlcolor=blue
}

\pagestyle{fancy}
\fancyhf{}
\lhead{Terminale spécialité -- Physique-Chimie}
\rhead{Corrigé détaillé du DM -- Électricité}
\cfoot{\thepage/\pageref{LastPage}}

\titleformat{\section}{\large\bfseries}{\thesection.}{0.5em}{}
\titleformat{\subsection}{\normalsize\bfseries}{\thesubsection.}{0.5em}{}

\setlist[itemize]{leftmargin=1.8em}
\setlist[enumerate]{leftmargin=2em}

\newcolumntype{Y}{>{\raggedright\arraybackslash}X}

\begin{document}

\begin{center}
    {\LARGE \textbf{Corrigé détaillé du Devoir Maison -- Physique-Chimie}}\\[0.3em]
    {\Large \textbf{Terminale spécialité}}\\[0.5em]
    {\large \textbf{Thème : Électricité}}\\[0.8em]
\end{center}



\section*{Problème 1 -- Charge d'un condensateur \hfill /24}

\[
E=\SI{12.0}{V},\qquad R=4{,}7\times10^3\,\si{\ohm},\qquad C=\SI{220}{\micro F}=2{,}20\times10^{-4}\,\si{F}
\]

\subsection*{1. Analyse qualitative du montage}

\begin{enumerate}
    \item Le circuit est un circuit série comportant un générateur idéal de tension, un interrupteur, une résistance \(R\) et un condensateur \(C\).

    \item Le courant conventionnel sort de la borne positive du générateur, traverse la résistance puis arrive sur l'armature du condensateur reliée au générateur.

    \item 
    \begin{enumerate}[label=\alph*)]
        \item À l'instant initial, le condensateur est déchargé, donc :
        \[
        u_C(0)=0.
        \]
        \item Au bout d'un temps très long, le condensateur est totalement chargé et :
        \[
        u_C(+\infty)=E.
        \]
        Le courant devient alors nul.
    \end{enumerate}

    \item Au cours de la charge, la tension \(u_C(t)\) augmente. Or la loi des mailles donne :
    \[
    u_R(t)=E-u_C(t).
    \]
    Ainsi, la tension aux bornes de la résistance diminue. Comme
    \[
    u_R(t)=Ri(t),
    \]
    l'intensité \(i(t)\) diminue au cours du temps.
\end{enumerate}

\subsection*{2. Équation différentielle}

\begin{enumerate}
    \item La loi des mailles donne :
    \[
    E=u_R(t)+u_C(t).
    \]

    \item Pour la résistance, la loi d'Ohm s'écrit :
    \[
    u_R(t)=Ri(t).
    \]

    \item Pour le condensateur :
    \[
    i(t)=C\frac{du_C}{dt}.
    \]

    \item En remplaçant \(u_R(t)\) puis \(i(t)\), on obtient :
    \[
    E=R\,i(t)+u_C(t)=RC\frac{du_C}{dt}+u_C(t).
    \]
    Donc
    \[
    RC\frac{du_C}{dt}+u_C=E.
    \]
\end{enumerate}

\subsection*{3. Solution et exploitation}

On admet :
\[
u_C(t)=E\left(1-e^{-t/RC}\right).
\]

\begin{enumerate}
    \item À \(t=0\),
    \[
    u_C(0)=E(1-e^0)=E(1-1)=0.
    \]
    La condition initiale est bien vérifiée.

    \item La constante de temps vaut :
    \[
    \tau=RC=4{,}7\times10^3\times2{,}20\times10^{-4}=1{,}034\ \si{s}.
    \]

    \item On a :
    \[
    u_C(\tau)=E(1-e^{-1}).
    \]
    Comme \(e^{-1}\approx 0{,}368\),
    \[
    u_C(\tau)\approx 12{,}0\times0{,}632\approx \SI{7.58}{V}.
    \]

    \item La constante de temps \(\tau\) est le temps caractéristique de la charge. Au bout d'une durée égale à \(\tau\), le condensateur a atteint environ \(63{,}2\,\%\) de sa tension finale.

    \item On dérive \(u_C(t)\) :
    \[
    \frac{du_C}{dt}=E\cdot \frac{1}{RC}e^{-t/RC}.
    \]
    Donc
    \[
    i(t)=C\frac{du_C}{dt}=C\cdot E\cdot \frac{1}{RC}e^{-t/RC}
    =\frac{E}{R}e^{-t/RC}.
    \]

    \item À \(t=0\),
    \[
    i(0)=\frac{E}{R}=\frac{12{,}0}{4{,}7\times10^3}
    \approx 2{,}55\times10^{-3}\ \si{A}
    =\SI{2.55}{mA}.
    \]

    \item Lorsque \(t\to+\infty\), on a \(e^{-t/RC}\to 0\). Donc
    \[
    i(t)\to 0.
    \]

    \item L'énergie stockée dans un condensateur vaut :
    \[
    E_C(t)=\frac12 C\,u_C^2(t).
    \]

    \item Lorsque le condensateur est totalement chargé, \(u_C=E=\SI{12.0}{V}\). Ainsi :
    \[
    E_C=\frac12\times 2{,}20\times10^{-4}\times 12^2.
    \]
    Or \(12^2=144\), donc
    \[
    E_C=\frac12\times 2{,}20\times10^{-4}\times 144
    =1{,}584\times10^{-2}\ \si{J}.
    \]
    Ainsi
    \[
    E_C\approx 1{,}58\times10^{-2}\ \si{J}.
    \]
\end{enumerate}

\subsection*{4. Regard critique}

\begin{enumerate}
    \item L'affirmation \og au début de la charge, le condensateur empêche totalement le passage du courant \fg\ est fausse.

    En effet, au début :
    \[
    u_C(0)=0,
    \]
    donc toute la tension du générateur est aux bornes de la résistance :
    \[
    u_R(0)=E.
    \]
    Ainsi le courant initial vaut
    \[
    i(0)=\frac{E}{R}\neq 0.
    \]
    Il est même maximal à cet instant.

    \item L'affirmation \og plus la résistance est grande, plus le condensateur se charge vite \fg\ est fausse.

    La vitesse caractéristique de charge est gouvernée par
    \[
    \tau=RC.
    \]
    Si \(R\) augmente, alors \(\tau\) augmente aussi, donc la charge est plus lente.

    \item Pour charger plus vite sans changer la tension finale atteinte, il faut diminuer la résistance \(R\). En effet, cela diminue \(\tau=RC\) sans modifier la valeur limite \(u_C(+\infty)=E\).
\end{enumerate}

\newpage

\section*{Problème 2 -- Capteur de présence \hfill /29}

\[
E=\SI{5.0}{V},\qquad R=\SI{100}{k\ohm}=1{,}00\times10^5\,\si{\ohm}
\]
\[
C_v=\SI{80}{nF}=8{,}0\times10^{-8}\,\si{F},\qquad
C_o=\SI{140}{nF}=1{,}40\times10^{-7}\,\si{F}
\]
\[
u_s=\SI{3.0}{V}
\]

\subsection*{1. Compréhension du principe}

\begin{enumerate}
    \item La capacité d'un condensateur dépend de la géométrie du système et du milieu environnant. La présence d'une personne modifie l'environnement électrique du capteur, ce qui modifie sa capacité.

    \item La charge est d'autant plus lente que la constante de temps
    \[
    \tau=RC
    \]
    est grande. Comme
    \[
    C_o>C_v,
    \]
    on a
    \[
    \tau_o>\tau_v.
    \]
    Le siège occupé correspond donc à une charge plus lente.

    \item La constante de temps d'un circuit RC vaut :
    \[
    \tau=RC.
    \]
\end{enumerate}

\subsection*{2. Temps de franchissement d'un seuil}

On rappelle :
\[
u_C(t)=E\left(1-e^{-t/RC}\right).
\]

\begin{enumerate}
    \item Cherchons \(t_s\) tel que \(u_C(t_s)=u_s\). On écrit :
    \[
    u_s=E\left(1-e^{-t_s/RC}\right).
    \]
    Donc
    \[
    \frac{u_s}{E}=1-e^{-t_s/RC},
    \]
    puis
    \[
    e^{-t_s/RC}=1-\frac{u_s}{E}.
    \]
    En prenant le logarithme népérien :
    \[
    -\frac{t_s}{RC}=\ln\left(1-\frac{u_s}{E}\right).
    \]
    D'où
    \[
    t_s=-RC\ln\left(1-\frac{u_s}{E}\right).
    \]

    \item Pour le siège vide :
    \[
    RC_v=1{,}00\times10^5\times8{,}0\times10^{-8}
    =8{,}0\times10^{-3}\ \si{s}
    =\SI{8.0}{ms}.
    \]
    Ensuite,
    \[
    1-\frac{u_s}{E}=1-\frac{3{,}0}{5{,}0}=0{,}4.
    \]
    Comme \(\ln(0{,}4)\approx -0{,}916\),
    \[
    t_{s,v}=-8{,}0\times10^{-3}\ln(0{,}4)
    =8{,}0\times10^{-3}\times0{,}916
    \approx 7{,}33\times10^{-3}\ \si{s}.
    \]
    Ainsi
    \[
    t_{s,v}\approx \SI{7.33}{ms}.
    \]

    \item Pour le siège occupé :
    \[
    RC_o=1{,}00\times10^5\times1{,}40\times10^{-7}
    =1{,}40\times10^{-2}\ \si{s}
    =\SI{14.0}{ms}.
    \]
    Donc
    \[
    t_{s,o}=-14{,}0\times10^{-3}\ln(0{,}4)
    =14{,}0\times10^{-3}\times0{,}916
    \approx 12{,}83\times10^{-3}\ \si{s}.
    \]
    Ainsi
    \[
    t_{s,o}\approx \SI{12.83}{ms}.
    \]

    \item On a
    \[
    t_{s,o}>t_{s,v}.
    \]
    Le siège occupé met plus de temps à atteindre le seuil, ce qui est logique car la capacité est plus grande et la charge plus lente.
\end{enumerate}

\subsection*{3. Détection automatique}

Le système décide que le siège est occupé si la tension \(u_C\) n'a pas encore atteint \(\SI{3.0}{V}\) au bout de \(\SI{9.0}{ms}\).

\begin{enumerate}
    \item 
    \begin{itemize}
        \item Cas du siège vide :
        \[
        t_{s,v}\approx \SI{7.33}{ms}<\SI{9.0}{ms}.
        \]
        Le seuil est donc atteint avant \(\SI{9.0}{ms}\).

        \item Cas du siège occupé :
        \[
        t_{s,o}\approx \SI{12.83}{ms}>\SI{9.0}{ms}.
        \]
        Le seuil n'est donc pas atteint à \(\SI{9.0}{ms}\).
    \end{itemize}

    \item La règle de décision est pertinente :
    \begin{itemize}
        \item siège vide : le seuil est atteint avant \(\SI{9.0}{ms}\), donc le système ne détecte pas une occupation ;
        \item siège occupé : le seuil n'est pas atteint à \(\SI{9.0}{ms}\), donc le système détecte correctement une occupation.
    \end{itemize}

    \item Toute valeur de temps comprise entre \(\SI{7.33}{ms}\) et \(\SI{12.83}{ms}\) conviendrait. Par exemple :
    \[
    t=\SI{10}{ms}
    \]
    serait un choix simple et robuste.
\end{enumerate}

\subsection*{4. Question de conception}

\begin{enumerate}
    \item 
    \begin{enumerate}[label=\alph*)]
        \item Si \(R\) augmente, alors \(\tau=RC\) augmente. Le temps de charge et donc le temps de mesure augmentent.
        \item Les temps caractéristiques étant proportionnels à \(R\), l'écart absolu entre les temps dans les cas \og vide \fg\ et \og occupé \fg\ augmente également.
    \end{enumerate}

    \item Une résistance trop grande rend le système lent : il devient moins réactif, ce qui peut être gênant pour une détection rapide.

    \item Une stratégie expérimentale consiste à tester plusieurs valeurs de \(R\), à mesurer les temps de réponse dans les deux situations, puis à choisir une valeur qui :
    \begin{itemize}
        \item sépare suffisamment bien les deux cas ;
        \item garde un temps de réponse raisonnable.
    \end{itemize}
\end{enumerate}

\subsection*{5. Raisonnement approfondi}

\begin{enumerate}
    \item À temps fixé,
    \[
    u_C(t)=E\left(1-e^{-t/RC}\right).
    \]
    Si \(RC\) est plus grand, alors la charge est moins avancée à l'instant considéré. Comme le siège occupé a une plus grande capacité, sa tension est plus faible à temps fixé. Donc la tension est plus grande pour le siège vide.

    \item Cette méthode peut être préférable car mesurer une tension à un instant donné est souvent plus simple électroniquement que mesurer précisément un temps de franchissement de seuil.

    \item Il faut choisir un instant \(t_m\) pour lequel les tensions des deux situations sont bien séparées. En pratique, on prend un temps intermédiaire, typiquement autour de \(\SI{9}{ms}\) ou \(\SI{10}{ms}\).
\end{enumerate}

\newpage

\section*{Problème 3 -- Déviation d'un électron dans un champ électrique uniforme \hfill /31}

\[
v_0=2{,}0\times10^7\,\si{m.s^{-1}},\qquad
E=1{,}5\times10^4\,\si{V.m^{-1}},\qquad
L=\SI{8.0}{cm}=8{,}0\times10^{-2}\,\si{m}
\]
\[
e=1{,}60\times10^{-19}\,\si{C},\qquad
m_e=9{,}11\times10^{-31}\,\si{kg}
\]

Le champ électrique est dirigé vers le bas.

\subsection*{1. Force subie par l'électron}

\begin{enumerate}
    \item La force électrique subie par une charge \(q\) placée dans un champ électrique \(\vec E\) est :
    \[
    \vec F=q\vec E.
    \]

    \item La charge de l'électron est négative :
    \[
    q=-e.
    \]

    \item Comme le champ électrique est dirigé vers le bas et que la charge de l'électron est négative, la force électrique est de sens opposé au champ. Elle est donc dirigée vers le haut.

    \item La valeur de la force est :
    \[
    F=eE=1{,}60\times10^{-19}\times1{,}5\times10^4
    =2{,}4\times10^{-15}\ \si{N}.
    \]
\end{enumerate}

\subsection*{2. Mouvement de l'électron}

\begin{enumerate}
    \item D'après la deuxième loi de Newton :
    \[
    m_e\vec a=\vec F.
    \]
    Donc
    \[
    \vec a=\frac{\vec F}{m_e}.
    \]
    Il n'y a pas de force selon \(Ox\), donc
    \[
    a_x=0.
    \]
    Selon \(Oy\),
    \[
    a_y=\frac{F}{m_e}
    =\frac{2{,}4\times10^{-15}}{9{,}11\times10^{-31}}
    \approx 2{,}63\times10^{15}\,\si{m.s^{-2}}.
    \]

    \item Comme \(a_x=0\), le mouvement selon \(Ox\) est rectiligne uniforme.

    \item Comme \(a_y\) est constant, le mouvement selon \(Oy\) est uniformément accéléré.

    \item Avec les conditions initiales
    \[
    x(0)=0,\quad y(0)=0,\quad v_x(0)=v_0,\quad v_y(0)=0,
    \]
    on obtient :
    \[
    x(t)=v_0t,
    \]
    \[
    y(t)=\frac12 a_y t^2.
    \]
\end{enumerate}

\subsection*{3. Trajectoire}

\begin{enumerate}
    \item Comme
    \[
    x=v_0t,
    \]
    on a
    \[
    t=\frac{x}{v_0}.
    \]

    \item En remplaçant dans l'expression de \(y(t)\), on obtient :
    \[
    y(x)=\frac12 a_y\left(\frac{x}{v_0}\right)^2
    =\frac{a_y}{2v_0^2}x^2.
    \]

    \item Cette équation est de la forme
    \[
    y=kx^2,
    \]
    avec \(k>0\). La trajectoire est donc une parabole.

    \item La durée de traversée est :
    \[
    t_s=\frac{L}{v_0}
    =\frac{8{,}0\times10^{-2}}{2{,}0\times10^7}
    =4{,}0\times10^{-9}\ \si{s}.
    \]

    \item La déviation verticale à la sortie vaut :
    \[
    y_s=\frac12 a_y t_s^2.
    \]
    Or
    \[
    t_s^2=(4{,}0\times10^{-9})^2=1{,}6\times10^{-17}.
    \]
    Donc
    \[
    y_s=\frac12\times2{,}63\times10^{15}\times1{,}6\times10^{-17}
    \approx 2{,}11\times10^{-2}\ \si{m}.
    \]
    Ainsi
    \[
    y_s\approx \SI{2.11}{cm}.
    \]
\end{enumerate}

\subsection*{4. Interprétation physique}

\begin{enumerate}
    \item La vitesse horizontale reste constante car aucune force n'agit selon \(Ox\). Donc \(a_x=0\).

    \item L'énergie cinétique varie. En effet, la vitesse verticale augmente sous l'action de la force électrique ; la norme de la vitesse augmente donc également. L'énergie cinétique
    \[
    E_c=\frac12 mv^2
    \]
    augmente au cours du mouvement.

    \item Si \(E\) augmente, la force électrique augmente en valeur absolue, donc l'accélération verticale augmente. La déviation verticale est donc plus grande.

    \item Si \(v_0\) augmente, le temps passé entre les plaques
    \[
    t_s=\frac{L}{v_0}
    \]
    diminue. L'électron a moins de temps pour être dévié, donc la déviation diminue.

    \item Augmenter \(E\) accentue la courbure de la trajectoire ; augmenter \(v_0\) la réduit.
\end{enumerate}

\subsection*{5. Ouverture}

\begin{enumerate}
    \item Sans modifier \(E\), on peut agir sur :
    \begin{itemize}
        \item la vitesse initiale \(v_0\),
        \item la longueur \(L\) des plaques.
    \end{itemize}

    \item Pour dévier davantage les électrons :
    \begin{itemize}
        \item il faut diminuer \(v_0\),
        \item il faut augmenter \(L\).
    \end{itemize}

    \item Une limite physique possible est qu'en augmentant trop \(L\) ou en diminuant trop \(v_0\), l'électron peut toucher une plaque avant la sortie.
\end{enumerate}


\end{document}